\section{Introduction}
\label{sec:introduction}

Heating, ventilation and air conditioning (HVAC) systems account for the
largest share of energy use in buildings, and buildings in turn account
for roughly a third of global final energy demand
\citep{Perez2008BuildingShare,IEA2024Buildings}. Decarbonising this end
use depends as much on better operation of existing systems as on
better hardware: advanced control strategies that exploit weather
forecasts, internal-gain signals, and slow envelope dynamics have
repeatedly been shown to deliver double-digit energy savings without
compromising occupant comfort
\citep{Drgona2020StochasticMPC,Mason2019RLReview}. Among these
strategies, deep reinforcement learning (DRL) has attracted particular
attention because it learns a closed-loop policy directly from
interaction with a building model, in principle without requiring
hand-tuned MPC weights or expensive solver invocations at run time.

A practical obstacle stands between the promise and the deployments,
however. Modern policy-gradient algorithms such as PPO consume on the
order of $10^{6}$--$10^{7}$ environment steps before they converge
\citep{Schulman2017PPO,Henderson2018Reproducibility}, and high-fidelity
building simulators are far too slow to provide those steps. The
BOPTEST testbed \citep{Blum2021BOPTEST}
addresses reproducibility by exposing a standardised HTTP API on top of
the IBPSA Modelica Buildings library \citep{Wetter2014Modelica}, but
on commodity hardware that API delivers fewer than two-dozen
environment steps per second --- two orders of magnitude below what a
single PPO training run requires per CPU thread. The standard
work-around in the literature is to replace the simulator with a
neural-network surrogate trained from logged trajectories
\citep{Zhang2019WholeBuilding,Amasyali2018Review,Kazmi2019DataDriven}.

When practitioners adopt the surrogate work-around they implicitly
adopt a stronger assumption as well: that a more physically faithful
surrogate must be a better RL training environment. The assumption is
intuitive, because a faithful surrogate keeps the policy's experience
close to what it would have been under the real simulator. It is also
under-tested in the literature. Most surrogate-vs.-RL studies report
either predictive validity (one-step or multi-horizon root-mean-square
error against held-out trajectories) or closed-loop controller
performance, not the relationship between the two
\citep{HouEvins2024,Bacher2011RCIdentification}. In
Section~\ref{sec:results_fidelity} we test the assumption explicitly
on BOPTEST \texttt{bestest\_air} and report a clear negative result.

The resolution we pursue keeps the calibrated physical twin in the
pipeline but moves it out of the rollout. The smoother data-driven
surrogate remains the training environment; the physical twin is
frozen after Stage A/B/C calibration and enters the policy training
loss only as a soft physical regulariser, penalising actions on which
the two surrogates physically disagree. This hybrid construction is
fast (Section~\ref{ssec:exp_runtime}) and faithful
(Section~\ref{sec:results_control}), and it exposes a non-trivial
empirical structure: the optimal regularisation strength is
\emph{controller-family specific} rather than universal.

\subsection*{Contributions}

This paper makes the following contributions:

\begin{enumerate}
    \item A physics-informed surrogate (\texttt{v3.5}) of the \boptest\
    \texttt{bestest\_air} testcase, calibrated through a three-stage inverse
    procedure (Stage A telemetry preprocessing, Stage B explicit identification
    of zone thermal capacitance $\czon$, Stage C residual head calibration),
    achieving a $38\%$ reduction in one-step temperature RMSE and a $40\%$
    reduction in power MAE over the uncalibrated baseline.
    \item A clear empirical demonstration that predictive validity does not
    transfer to reinforcement learning training utility: the calibrated twin
    achieves a 24-hour rollout RMSE of \SI{0.64}{\celsius} on held-out
    trajectories yet drives a closed-loop \boptest{} RMSE above
    \SI{4}{\celsius} when used as a stand-alone RL training environment.
    \item A hybrid surrogate construction (\texttt{hybrid\_l010}) in which the
    calibrated twin acts as a frozen \emph{soft physical regularizer} on a
    smoother control-oriented surrogate, restoring closed-loop transfer RMSE
    below \SI{0.65}{\celsius} on both \texttt{peak\_heat\_window} and
    \texttt{typical\_heat\_window} scenarios while sustaining
    \num{1787} environment steps per second on a single CPU thread
    ($85.0\times$ speed-up over the live \boptest{} RTE HTTP loop at the
    same \SI{15}{\minute} control protocol; see
    Table~\ref{tab:speed}).
    \item A controller-family-specific finding: the optimal physical
    regularization strength differs across controller architectures
    ($\lambda_{\mathrm{temp}}=0.10$ for thermostatic PPO,
    $\lambda_{\mathrm{temp}}=0.00$ for HDRL and 17-D MORL), with an
    interpretation grounded in observation-space geometry and control
    hierarchy.
    \item A pre-registered transferability characterisation across three
    additional \boptest{} hydronic testcases that decomposes the
    recipe into a surrogate component, which transfers structurally
    (identified $\czon$ ratio $1.91\pm0.03$ on $N=3$), and a
    controller component, which fails on all three under a frozen
    direct-supply-temperature policy. The decomposition is the central
    Block~3 result and is reported in Section~\ref{sec:results_transfer}.
    \item A reproducible \boptest-based evaluation protocol, with full
    \cite{HouEvins2024}-compliant numerical justification for sample
    generation, preprocessing, feature significance and independence,
    scaling, hyperparameters, architecture, and replicative-plus-predictive
    validity, provided as Supplementary Tables S1--S11.
\end{enumerate}

The remainder of the paper is organised as follows.
Section~\ref{sec:related_work} positions the work against three
literatures (DRL for HVAC, surrogate modelling for buildings, and
physics-informed machine learning).
Section~\ref{sec:methodology} describes the two surrogates, the
hybrid loss, and the controller families.
Section~\ref{sec:experimental_setup} fixes the evaluation protocol,
introduces the \cite{HouEvins2024} compliance and the
pre-registration audit chain.
Sections~\ref{sec:results_fidelity}--\ref{sec:results_transfer}
report the three blocks of results in order: surrogate fidelity,
controller performance, and transferability across testcases.
Section~\ref{sec:discussion} explains the mechanisms behind the
controller-family-specific findings and states the threats to
validity, and Section~\ref{sec:conclusion} closes with a tight scope
statement and the natural follow-up.
